% Source-audit fragment, 2026-10-07. No documentclass; not inserted in main.
% This identifies missing interfaces and gives valid replacement formulas.
% It neither refutes Li's theorem nor certifies its numerical constant.

\subsection{What the cited double-sieve bounds actually transfer}
\label{sec:external-chen-input-resolution}

Let $d=N^{T+o(1)}$ be a specified ordered prime prefix and let
$z=N^u$. Write
\[
 s_0=\frac{1/2-T}{u},\qquad s_1=\frac{\nu-T}{u},
 \qquad \nu>1/2.
\]
All statements below concern an averaged positive prefix box,
with the common main factor denoted by $M>0$. They do not assert
an estimate for each individual prefix prime.
The improved linear estimate at level $N^\nu$ must itself be
justified on that box.

In Wu's \href{https://arxiv.org/pdf/0709.3764}{2008 author version},
Section 3, the total level is $Q=N^{1/2-\delta}$, the residual
level is $Q/d$, and $A(s)=sF(s)/(2e^\gamma)$,
$a(s)=sf(s)/(2e^\gamma)$.
The functions $H,h$ are defined through the defect from these
same-level normalized functions in (3.4)--(3.7).
Their prefix class satisfies
$p_1\cdots p_{j-1}p_j^2\le Q$.
This is also the normalization underlying Section 3 and Lemmas
5.1--5.2 of
\href{https://www.impan.pl/shop/en/publication/transaction/download/product/82961}
{Wu's 2004 paper}; the latter's functional inequalities iterate
the same total level, rather than adding an independent saving
to a different linear-sieve baseline.

Thus, after removing fixed margins, the two separately established
upper estimates have the form
\[
 S/M\le F(s_1)+o(1),\qquad
 S/M\le F(s_0)-D_H+o(1),\qquad
 D_H=\frac{2e^\gamma}{s_0}H(s_0).
\]
The analogous two lower estimates use $f(s_1)$ and
$f(s_0)+D_h$, where $D_h=2e^\gamma h(s_0)/s_0$.
Whenever both inputs are valid on the same positive box, their
rigorous combination is
\begin{align}
 U_{\rm safe}
 &=\min\{F(s_1),F(s_0)-D_H\}\notag\\
 &=F(s_1)-[D_H-\{F(s_0)-F(s_1)\}]_+,
 \label{eq:external-safe-upper}\\
 L_{\rm safe}
 &=\max\{f(s_1),f(s_0)+D_h\}\notag\\
 &=f(s_1)+[D_h-\{f(s_1)-f(s_0)\}]_+.
 \label{eq:external-safe-lower}
\end{align}
Here $[a]_+=\max(a,0)$. One can use certified lower bounds for
$H,h$ in these formulas. A finite common partition permits the
choice of the better bound on each box before taking the
large-$N$ limit. On boxes outside Wu's prefix class, omit the
classical double-sieve option.

These formulas follow by taking a minimum or maximum of bounds
for the same quantity. They do not require a new double-sieve
theorem. Relative to the mixed expressions
$U_{\rm mix}=F(s_1)-D_H$ and $L_{\rm mix}=f(s_1)+D_h$, the exact
local corrections needed are
\begin{align*}
 U_{\rm safe}-U_{\rm mix}
   &=\min\{D_H,F(s_0)-F(s_1)\},\\
 L_{\rm mix}-L_{\rm safe}
   &=\min\{D_h,f(s_1)-f(s_0)\}.
\end{align*}
They quantify the loss that a corrected integral calculation
must include. Merely imposing the two separate bounds cannot
remove that loss: for example, when
$F(s_1)\le F(s_0)-D_H$ and $D_H>0$, the hypothetical value
$S/M=F(s_1)$ satisfies both bounds but violates $U_{\rm mix}$.
This is a logical non-implication, not a counterexample to a
prime-counting theorem.

\paragraph{The additional statement needed for the full mixed saving.}
Lichtman,
\href{https://arxiv.org/html/2109.02851v2}{arXiv:2109.02851v2},
Lemma 6.4 and (6.23), give the upper correction
\[
 F(s_1)-\frac{2e^\gamma}{s_1}H_\nu(s_1),
 \qquad s_1\in[2,3],
\]
subject to distribution at the same level for every extended
prime vector and the corresponding switched sets.
The largest specified prime is preserved during the iteration
in the application following that lemma. The stated monotonicity
is $H_\nu(s)\ge H_{1/2}(s)$ at a fixed argument $s$.
To deduce the full old-argument saving used in Li's formulas
from this statement one would additionally need
\begin{equation}\label{eq:external-missing-cross-level}
 H_\nu(s_1)\ge\frac{s_1}{s_0}H_{1/2}(s_0),
 \qquad
 h_\nu(s_1)\ge\frac{s_1}{s_0}h_{1/2}(s_0),
\end{equation}
or a direct joint sieve construction proving the desired mixed
bounds. Fixed-argument monotonicity does not imply these
inequalities; $H_{1/2}$ decreases as its argument increases.
For the lower correction, the final paragraph of Lichtman's
Section 6.3 explicitly leaves the analogous formula derivations
to the reader. Extension of its stated $H$ lemma beyond $[2,3]$
also requires an argument at the new level.

The checked passages in
\href{https://imar.ro/journals/Mathematical_Reports/Pdfs/2026/1_2/3.pdf}
{Li II}, (16), (19), (21), and (23)--(25), use the old arguments
$(1/2-T)/u$ in $H,h$ alongside improved linear terms. The
paragraph following (16) cites fixed-argument monotonicity.
The earlier \href{https://arxiv.org/html/2405.05727v3}{2024 v3}
already has the same combination and does not supply a further
cross-level lemma. Consequently the checked citations do not
close \eqref{eq:external-missing-cross-level}. This audit has
not proved those inequalities false.

\paragraph{The classical prefix domains are consistent.}
Put $a=100/1149$, $h=50/309$, $\mu=25/128$, and
$d_*=1/2-3a$. The old double-sieve domains in Li (23)--(25)
are compatible with the old support class. For (23),
$t_1+2t_2\le3h<1/2$. In (24), the split at $1/2-2h$
and the later bound $t_2\le39/256=(1/2-\mu)/2$ ensure
$t_1+2t_2\le1/2$. In the correction of (25),
$t_1\le d_*$ and $t_2\le3a/2$ give the same inequality.
Also $2t_1<1/2$ throughout these regions.
Thus the old-level terms are usable as separate bounds;
the issue is their addition to the new baseline.
Where only tabulated values for $s\ge2$ are supplied, using
$h(s)\ge0$ for smaller $s$ is a valid conservative choice.

\subsection{The largest-prime restriction is active}
\label{sec:external-largest-prime-resolution}

Let $\Phi$ denote the raw eight-branch maximum of L23 (6.4),
and let $v(t)$ be its single-variable level. Li's Section 2
refers to these functions and states the three-variable
improvement only when the largest prime exponent is at most
$\mu$. No piecewise extension of the three-variable function
was located in the checked Li II text or its cited 2024 v3.
In contrast, L23 itself uses the three-variable expression in
(7.5) only for $c\le\mu$ and explicitly returns to $v(t)$
in (7.6) for $c>\mu$.

A conservative level supported by the stated distribution
theorems, with their usual strict margins, is therefore
\[
 \Phi_{\rm safe}(x,y,z)=
 \begin{cases}
  \Phi(x,y,z),&x\le\mu,\\
  v(x),&x>\mu,
 \end{cases}
 \qquad x\ge y\ge z.
\]
The full prefix interface is still required wherever the first
case is invoked; this notation does not bypass its hypotheses.
For (34), membership in $D^{\rm well}$ must also be recomputed
at the selected safe level; when it fails, retain only the
switched estimate.

This restriction is not rendered moot by vanishing integrands.
In (25), take
\[
 (t_1,t_2,u)=(11/50,3/20,100/1149).
\]
This lies strictly inside its two-dimensional integration
domain. Here
\[
 \Phi=\frac{23}{40},\quad v(t_1)=\frac{897}{1600},\quad
 s_{\rm raw}=\frac{47109}{20000}>2,\quad
 s_{\rm safe}=\frac{70089}{32000}>2.
\]
Since both arguments lie between 2 and 4, the lower-sieve
function is strictly increasing there. Hence the raw and safe
lower terms differ on an open neighborhood.

For (34), take the ascending tuple
\[
 (t_1,t_2,t_3,t_4)=(11/125,9/100,23/250,101/500).
\]
It is strictly inside the displayed integration domain, with
$t_4>\mu$. The full four-prime square-prefix exponents are
\[
 101/250,\quad193/500,\quad237/500,\quad14/25.
\]
They are all less than both
\[
 \Phi(t_4,t_3,t_2)=\frac{591}{1000},\qquad
 v(t_4)=\frac{4701}{8000}.
\]
The raw maximizing pair branch has gate surplus $1/200$.
Writing $S=59/125$, the two upper-sieve arguments are
$119/88$ and $925/704$, respectively. Both exceed 1, so neither
the support indicator nor a zero-density convention removes
this example. With the published switched coefficient cap
$G_1=6.06932$, the linear option is also smaller than the
switched option here: their ratio is
\[
 \frac{F((\nu-S)/t_1)t_2}
 {G_1\omega((1-S)/t_2)t_1}
 =\frac{2e^\gamma t_2}{G_1\omega((1-S)/t_2)(\nu-S)}<1
\]
for either displayed $\nu$, using
$\omega(u)\ge1/2$ for $u\ge2$, $e^\gamma<1.782$, and
$\nu-S\ge37/320$. Thus the minimum need not hide the issue.

The domain restriction has a concrete repair and explicit
nonzero affected regions. The surviving task is to recompute
the affected functionals, with
\eqref{eq:external-safe-upper}--\eqref{eq:external-safe-lower}
or a proved stronger joint construction. No new numerical
Chen coefficient is certified in this fragment. In particular,
these findings neither independently establish nor disprove
the published value $1.9728$.
